\section{Extended Limitations}

SPINE is a transparent computational framework, not a calibrated biological reconstruction. The baseline reference track preserves pre-specified assumptions even when they are simplified. The plausibility-revised and active-extension tracks are explicitly separated. Exploratory perturbation scenarios are supplemental stress tests and are not disease models.

The analytic divider is a DC small-signal reference. The simulated targets in this manuscript are transient conductance-based peak responses, so residuals from the divider are expected and informative. The independent direct-matrix benchmark checks the passive three-compartment baseline implementation against an independently assembled equation solve, but it does not check every morphology, active, uncertainty, or exploratory output. NEURON cross-checking was not performed.

The deterministic uncertainty ensembles are sensitivity designs, not experimentally fitted posterior distributions. The N=768 ensemble contains no high-SMI rows, so high-isolation uncertainty remains a future targeted extension. Radius uncertainty, unmodeled ultrastructure, cytoplasmic resistivity uncertainty, segmentation error, local synaptic receptor composition, and active conductance distributions can all alter interpretation.

The current package does not include posterior inference from experimental voltage traces, cell-type-specific channel fitting, full morphology libraries, biochemical calcium reaction-diffusion, receptor trafficking, signaling cascades, or public release approvals. These are future validation and administrative targets, not hidden assumptions of the current manuscript.
